\documentclass[
    a4paper,
    DIV=13,
    abstract=true,
    numbers=noenddot
]{scrartcl}

\usepackage{preamble}
\addbibresource{library.bib} % duplicate of preamble, tells VimTeX where entries live

\begin{document}
\title{WGF Intro}
\author{Liam Llamazares-Elias}
\date{\today}
\maketitle
\section{Introduction}

\section{Euclidean gradient flows}

\section{Convergence}
Gronwall.
PL theorem. CIte Pablo.
Convergece to subspace theorem

\begin{theorem}[PL]\label{thm:PL}
Let $\Phi: \R^d \to \R$ be locally Lipschitz and verify the PL inequality 
\begin{aligneq*}
    \norm{\nabla \Phi(x)}^2 \ge \lambda \norm{\Phi(x)},\quad \forall x \in \R^d. 
\end{aligneq*}
Then, a solution $x_t$ to the gradient flow 
\begin{aligneq*}
   x_t'=-\nabla \Phi(x_t), \quad x_0=x_0
\end{aligneq*}
has a unique stationary point $x^*(x_0)$ and verifies the exponential convergence 
\begin{aligneq*}
    \norm{x_t-x^*(x_0)}^2 \le \exp(- \lambda t)\norm{x_t-x^*(x_0)}.
\end{aligneq*}
\end{theorem}
\begin{proof}
    By the assumption of the theorem a unique solution exists. Using the PL ineuqality 
\begin{aligneq*}
    \partial_t \Phi(x_t)= - \norm{\nabla \Phi(x_t)}^2 \le -\lambda \norm{\Phi(x_t)}.
\end{aligneq*}
As a result, by Gronwall's inequality, $\Phi$ decreases expoentially 
\begin{aligneq*}
  \Phi(x_t) \le e^{-\lambda t}\Phi(x_0).
\end{aligneq*}
Since $\Phi$ is the only thing moving $x$, the trajectory slows down and has finite length 
\begin{aligneq*}
  \int_{0}^\infty \norm{x'_t} \d t =\int_{0}^\infty \norm{\Phi(x_t)} \d t \le=\int_{0}^\infty e^{-\lambda t}\norm{\Phi(x_0)} \d t < \infty.      
\end{aligneq*}
A small argument confirms that any curve of finite length is Cauchy and thus converges to a unique point.
\end{proof}
A conventional example is that of strongly convex functions?
\begin{exercise}
 Show that if $\Phi$ has   
    \exerciseSymbol
\end{exercise}
A typical example occurs when $Q$ is positive definite and $\Phi=\frac{1}{2} \br{ Qx,x }$ is a quadratic potential. In which case it is clear that the constant $\lambda$ governing the speed of convergence is the smallest eigenvalue of $Q$.

\section{Preconditioning}
When $\lambda \approx 0$ the speed of convergence is slow as  we require time $t= \lambda^{-1} \log{\eps}$ to be within a distance $\eps$ of $x^*$. This can be improved through preconditioning. That is, by sugin a second order method. Consider the previous case $\Phi= \br{ Qx,x }$ where $ \lambda \sim 0$. Weuse $\Sigma:= Q ^{-1}$ to straighten out the geometry, considering 
\begin{aligneq*}
    y'_t= -\Sigma \nabla \Phi(y_t),
\end{aligneq*}
Now, by \Cref{thm:PL}, 
\begin{aligneq*}
  \norm{x_t -x^*} \le \norm{Q}\norm{y_t-y_0} \le \lambda_{\rm{max}} e^{-t}.
\end{aligneq*}
So achieving an error of $\eps$ requires $t= \log{\lambda_{\rm{max}} \eps^{-1}}$ instead of \liam{Surepsed amx eigen comes up, what if it is very big. Then we are worse}


\printbibliography
\end{document}
